% Gerado por ferramentas/importar_relatorios.py a partir de CWRU/RELATORIO_CWRU-10.11.2026/RELATORIO_CWRU.tex.
% Revisado à mão na dissertação: edite este arquivo, não a fonte.
\subsection{Objetivo e enquadramento}
\label{cwru:sec:objetivo}

O estudo da bomba (Seção~\ref{sec:mm_bomba}) mostrou que modelos treinados com descritores de \textit{hits} reconhecem bem as condições quando a corrida aparece no treino, mas generalizam mal quando a corrida inteira é deixada de fora, e que os dados tinham problemas de aquisição: 70\,\% dos \textit{hits} saturados e muitos truncados pela duração máxima. Como os dados do ensaio têm uma única corrida por condição, não era possível separar o que é limitação do método do que é limitação dos dados.

A base CWRU \cite{cwru} é a referência mais usada em diagnóstico de rolamentos e permite testar o mesmo fluxo em dados públicos e rotulados. Ela tem, porém, uma estrutura parecida com a do ensaio da bomba: cada rolamento com defeito foi montado uma vez e gravado em quatro cargas, de modo que há poucas unidades físicas independentes. Parte da literatura relata acurácias próximas de 100\,\% com janelas sorteadas, e trabalhos posteriores mostraram que esses números caem quando a divisão respeita as gravações ou os rolamentos \cite{smith2015,rauber2021,hendriks2022}.

Os sinais são de vibração (acelerômetros a 48\,kHz), e não de EA. O extrator de \textit{hits} do MOMEA é aplicado aqui como detector de impulsos: os defeitos localizados produzem impactos periódicos, e a lógica de limiar, DDT e rearme se aplica a eles. Os descritores, porém, têm outro significado físico, e as conclusões sobre cavitação não se transferem. O objetivo é responder a três perguntas:
\begin{enumerate}
	\item os descritores de \textit{hits} no estilo MOMEA capturam a física do defeito?
	\item como a acurácia depende da validação, quando as unidades físicas são separadas?
	\item o fluxo embarcado (MLP int8) preserva a decisão?
\end{enumerate}

% =======================================================================================================
